\documentclass[11pt]{article}
\usepackage[T1]{fontenc}
\usepackage{lmodern,amsmath,amssymb,amsthm,booktabs}
\usepackage[letterpaper,margin=1in]{geometry}
\usepackage{microtype,xcolor}
\usepackage[colorlinks=true,linkcolor=blue!45!black,urlcolor=blue!45!black]{hyperref}
\newtheorem{theorem}{Theorem}[section]
\newtheorem{lemma}[theorem]{Lemma}
\newtheorem{proposition}[theorem]{Proposition}
\newtheorem{corollary}[theorem]{Corollary}
\newtheorem{remark}[theorem]{Remark}
\setlength{\emergencystretch}{3em}
\title{Exponential small-cut estimates for uniform even subgraphs}
\author{Research working draft}
\date{10 November 2026}
\begin{document}
\maketitle
\begin{abstract}
Let $U_1,\ldots,U_k$ be disjoint connected cyclic regions in a connected graph.
Suppose that each region has at most $d\ge2$ boundary edges.
For a uniform even subgraph, we prove that the probability of a forest in every region is at most
\[
 \frac12+2\exp\!\left(-\frac{k}{8d2^d}\right).
\]
The proof combines local coordinate spaces, vertex splitting, terminal forests, and conditional independence across blocks.
A weighted version uses the sum of the boundary-closure probabilities.
Parallel bundles on a cycle show that the degree scale $2^{-d}$ is necessary up to a factor of order $d$.
The proofs have passed independent mathematical audits, including an audit by Fable 5.2.
They are not formalized in Lean, and publication priority has not been established.
\end{abstract}

\section*{Reading guide}
Rank means $|E|-|V|+\operatorname{comp}$.
An even subgraph is an edge set with even degree at every vertex.
An internal even subgraph is a forest exactly when it is empty.
Parallel edges remain distinct.

The first section proves the common-functional mechanism and a coarse exponential estimate.
The second section uses graphical structure to improve its constant.
The third section proves the bound stated in the abstract.
The last two sections give the weighted form and the degree-dependence obstruction.
Intermediate estimates are retained because their mechanisms may be useful in further work.

\tableofcontents
% New research fragment. Root and path_stability independently audited the safe exponential bound.
% Do not merge into the previously audited manuscripts without a separate review.
\section{Exponential decay for disjoint regions with bounded cuts}
\label{sec:second-order-exponential-cuts}

All cycle spaces below are over $\mathbb F_2$.
Graphs may have parallel edges.
For a graph $Q$, write $\mathcal Z(Q)$ for its binary cycle space.
A uniform element of this space is denoted by $F$.
A vertex is \emph{closed} when no incident edge belongs to $F$.
Loops can be deleted when discussing this event, since it concerns the edge boundary of a vertex.

\begin{lemma}[A common local functional]\label{lem:common-star-functional}
Let $Q$ be a connected loopless multigraph.
Let $v_1,\ldots,v_\ell$ be distinct vertices of degree at most $d$, where $d\ge2$.
Suppose a nonzero functional $a\in\mathcal Z(Q)^*$ lies in the span of the incident edge-coordinate functionals at every $v_i$.
If $Z$ counts the closed vertices among $v_1,\ldots,v_\ell$, then
\[
 \mathbb E\,2^{-Z}
 \le\frac12+\frac12\left(1-2^{1-d}\right)^{\lfloor\ell/2\rfloor}.
\]
\end{lemma}

\begin{proof}
At $v_i$, choose an incident-edge subset $A_i$ whose coordinate sum represents $a$.
It is nonempty and proper: the empty subset and the whole star both represent zero.
Split $v_i$ into two vertices carrying $A_i$ and its complement, and join them by a new edge $h_i$.
Perform all splits simultaneously, obtaining a connected graph $Q'$.
Contracting the new edges identifies the two cycle spaces.
Under this identification, the coordinate on every $h_i$ equals $a$.

Deleting all $\ell$ new edges imposes exactly the one condition $a=0$.
Thus
\[
 \beta\bigl(Q'-\{h_1,\ldots,h_\ell\}\bigr)=\beta(Q')-1.
\]
Keeping isolated vertices in the rank calculation shows that the remaining graph has exactly $\ell$ components.
Contract these components.
The resulting connected quotient has $\ell$ vertices and $\ell$ edges.
An even subgraph with $a=1$ projects to its entire edge set.
The quotient is therefore connected and two-regular, hence a cycle.
For $\ell=1$ it is a loop, and for $\ell=2$ it consists of two parallel edges.

Condition on $a=0$.
The restrictions to the $\ell$ components are independent uniform even subgraphs.
Choose a matching of $\lfloor\ell/2\rfloor$ edges in the quotient cycle.
For each chosen $h_i$, its two split endpoints belong to distinct components.
Different chosen edges use disjoint component pairs.

Suppose the two parts of the original star have sizes $s_i$ and $t_i$.
Their sum is $\deg_Q(v_i)\le d$.
In its component, a vertex of degree $s_i$ has at most $s_i-1$ independent incident-edge coordinates.
Its star is therefore zero with probability at least $2^{1-s_i}$.
The corresponding bound at the other endpoint is $2^{1-t_i}$.
Consequently, the original vertex $v_i$ is closed with conditional probability at least $2^{2-d}$.
For the chosen matching, these events are independent.
It follows that
\[
 \mathbb E[2^{-Z}\mid a=0]
 \le\left(1-2^{1-d}\right)^{\lfloor\ell/2\rfloor}.
\]
When $a=1$, none of the selected original vertices is closed.
Since $a$ is a nonzero functional, its two values have equal probability.
This proves the result.
For $\ell=1$, the displayed bound is the trivial bound one.
\end{proof}

\begin{theorem}[Exponential closed-star estimate]\label{thm:exponential-closed-stars}
Let $Q$ be a connected loopless multigraph with $k\ge1$ distinct marked vertices of degree at most $d$.
Let $Z$ count the marked vertices closed in a uniform even subgraph.
For $d\ge2$, put
\[
 T_d=2^d,\qquad D_d=(d-1)(2^d-1),\qquad M_d=2^{D_d}-1.
\]
Then
\begin{equation}\label{eq:exponential-closed-stars}
 \mathbb E\,2^{-Z}
 \le\frac12+\exp\!\left(-\frac{k}{2^d M_d}\right)
 \le\frac12+\exp\!\left(-\frac{k}{2^{D_d+d}}\right).
\end{equation}
For $d\le1$, every marked vertex is closed and $\mathbb E2^{-Z}=2^{-k}$.
\end{theorem}

\begin{proof}
Let $L_i\subseteq\mathcal Z(Q)^*$ be the span of the edge-coordinate functionals at marked vertex $i$.
The parity equation at that vertex gives $\dim L_i\le d-1$.
If some $L_i$ is zero, that vertex is always closed.
The expectation is then at most one half, so assume every $L_i$ is nonzero.

If $T_d$ of these spaces form a direct sum, their local patterns are independent.
Each corresponding vertex is closed with probability at least $2^{1-d}$.
Consequently,
\[
 \mathbb E2^{-Z}\le(1-2^{-d})^{T_d}\le e^{-1}<\frac12.
\]

Otherwise choose a maximal direct-sum family and let $U$ be its span.
It has fewer than $T_d$ members, so $\dim U\le D_d$.
Every unselected $L_i$ meets $U$ nontrivially, by maximality.
Every selected $L_i$ also meets $U$ nontrivially, since it is nonzero and contained in $U$.
There are at most $M_d$ nonzero vectors in $U$.
Thus one nonzero functional $a$ belongs to at least $\ell\ge k/M_d$ of the spaces $L_i$.

Apply Lemma~\ref{lem:common-star-functional} to those vertices and ignore the other marked factors.
With $q=1-2^{1-d}\ge1/2$, one has
\[
 \frac12q^{\lfloor\ell/2\rfloor}
 \le\frac1{2q}\exp\!\left(-\frac{(1-q)\ell}{2}\right)
 \le\exp(-2^{-d}\ell).
\]
This proves~\eqref{eq:exponential-closed-stars}.
The assertion for $d\le1$ follows directly from parity.
\end{proof}

\begin{corollary}[Disjoint cyclic regions]\label{cor:exponential-small-cuts}
Let $G$ be a connected graph and let $U_1,\ldots,U_k$ be pairwise disjoint vertex sets.
Suppose each $G[U_i]$ is connected and contains a cycle, and $|\delta_G(U_i)|\le d$.
For a uniform even subgraph $F$ of $G$, let $\mathcal F$ be the event that every $F\cap E(G[U_i])$ is a forest.
For $d\ge2$,
\[
 \Pr(\mathcal F)
 \le\frac12+\exp\!\left(-\frac{k}{2^d(2^{(d-1)(2^d-1)}-1)}\right).
\]
For $d\le1$, one has $\Pr(\mathcal F)\le2^{-k}$.
\end{corollary}

\begin{proof}
Contract each connected induced region to a marked vertex and delete the resulting internal loops.
Retain all other vertices and edges, obtaining a connected quotient $Q$.
Projection maps $\mathcal Z(G)$ onto $\mathcal Z(Q)$.
Its kernel is the direct sum of the cycle spaces of the induced regions.
Thus the quotient pattern is uniform.
Conditional on that pattern, the internal choices in different regions are independent affine cycle-space choices.

If a marked quotient vertex is closed, the corresponding internal choice is a uniform even subgraph of $G[U_i]$.
Its only forest is the empty edge set.
Since that region has positive rank, its forest probability is at most one half.
For a nonclosed marked vertex, bound the conditional forest probability by one.
If $Z$ is the number of closed marked vertices, these observations give
\[
 \Pr(\mathcal F\mid\text{quotient pattern})\le2^{-Z}.
\]
Any remaining quotient loops are at unmarked vertices.
Delete them before applying the loopless closed-star theorem.
Their independent coordinates do not change any marked closure event.
Average and apply Theorem~\ref{thm:exponential-closed-stars}.
\end{proof}

\paragraph{Quantitative scope.}
For each fixed cut bound $d$, the error decreases exponentially with the number of disjoint cyclic regions.
The explicit constant is large: its binary logarithm is $O(d2^d)$.
This theorem improves the previous inverse-linear dependence on the region count.
It does not by itself establish a polynomial error for unrestricted cycle counts.
When applying it with a growing cut bound, the dependence on $d$ must be retained.

% Independent research fragment. Root derivation; independent audits passed.
\section{A terminal-forest improvement to the cut constant}
\label{sec:terminal-forest-cut-constant}

The common-functional lemma gives a much better constant when the local spaces come from a graph.
The following argument replaces a vector-space pigeonhole by a forest with few terminals.

\begin{lemma}[A local intersection in a terminal forest]
\label{lem:local-terminal-intersection}
Let $Q$ be a connected loopless multigraph, and let $S$ be a nonempty proper vertex set.
Write $L_v$ for the span of the edge-coordinate functionals at $v$ in $\mathcal Z(Q)^*$.
Put $U=\sum_{v\in S}L_v$ and $B=|\delta_Q(S)|$.
Let $\mathcal V$ be a set of vertices outside $S$ with $L_v\cap U\ne0$.
Except for at most $2B$ members, $\mathcal V$ can be partitioned into at most $2B$ classes.
Each class has a common nonzero functional in all its spaces $L_v$.
\end{lemma}

\begin{proof}
At most $B$ vertices outside $S$ have a neighbor in $S$.
Exclude those vertices.
Choose a spanning tree in each component of $Q-S$.
For each edge of $\delta_Q(S)$, attach a distinct terminal leaf to its endpoint outside $S$.
Let $T$ be the union of paths between terminal leaves in these trees.
Components with fewer than two terminals contribute nothing.
The forest $T$ has at most $B$ leaves.
It has at most $B$ vertices of degree at least three.
After suppressing its degree-two vertices, it has at most $2B$ edges.
Exclude the members of $\mathcal V$ at branching vertices of $T$.

We show that every remaining member lies internally on one of these edges.
Fix $v\in\mathcal V$ with no neighbor in $S$, and choose $0\ne a\in L_v\cap U$.
Represent $a$ both by a subset $A$ of the star at $v$ and by a subset of edges incident with $S$.
The symmetric difference of these two representing sets annihilates $\mathcal Z(Q)$.
It is therefore a cut of $Q$.
Let $H$ be the component of $Q-S$ containing $v$.
Every component of $H-v$ lies entirely on one side of this cut.
Consequently, $A$ consists of the edges from $v$ to a union of components of $H-v$.
As a functional, $a$ equals the sum of the boundary-terminal coordinates attached to that union.

The sum of all terminal coordinates of $H$ is zero, because $\delta_Q(H)$ is a cut.
Since $a\ne0$, at least two components of $H-v$ contain terminals.
Every spanning tree of $H$ must connect these terminal groups through $v$.
Thus $v$ has degree at least two in $T$.
It has degree exactly two after the branching exceptions have been removed.
There are exactly two terminal-bearing components of $H-v$.
Their terminal sets are the two sides of the corresponding edge of the suppressed forest.
The sum of the terminal coordinates on either side is $a$.
In particular, this common functional is nonzero.

All remaining members on the same suppressed edge give the same terminal bipartition.
They therefore contain the same nonzero functional in their local spaces.
The adjacent and branching exceptions number at most $2B$.
The number of classes is at most $2B$, as required.
\end{proof}

\begin{theorem}[An explicit exponential cut estimate]
\label{thm:terminal-forest-cut-estimate}
Let $Q$ be a connected loopless multigraph with $k$ distinct marked vertices of degree at most $d\ge2$.
Let $Z$ count the marked vertices closed in a uniform even subgraph.
Then
\[
 \mathbb E2^{-Z}\le\frac12+2\exp\!\left(-\frac{k}{2d4^d}\right).
\]
The same bound applies to the probability that a uniform even subgraph is a forest in each of $k$ disjoint cyclic regions.
Each region must have an edge boundary of size at most $d$.
\end{theorem}

\begin{proof}
Use the local spaces $L_v$ from Theorem~\ref{thm:exponential-closed-stars}.
A zero local space gives an expectation at most one half.
If $2^d$ local spaces form a direct sum, that theorem's independence argument also gives an expectation below one half.

Otherwise take a maximal direct-sum family, with vertex set $S$ and size $t<2^d$.
Every remaining local space meets $U=\sum_{v\in S}L_v$ nontrivially.
Suppose first that $k>(1+2d)2^d$.
Then $S$ is nonempty and proper, and $B=|\delta_Q(S)|\le dt$.
Lemma~\ref{lem:local-terminal-intersection} gives a common nonzero functional for at least
\[
 \ell\ge\frac{k-t-2B}{2B}
 \ge\frac{k}{2d2^d}-2
\]
marked vertices.
The common-functional lemma, with $q=1-2^{1-d}\ge1/2$, gives
\[
 \mathbb E2^{-Z}-\frac12
 \le\frac12q^{\lfloor\ell/2\rfloor}
 \le e^{-2^{-d}\ell}
 \le e^{2^{1-d}}\exp\!\left(-\frac{k}{2d4^d}\right)
 \le2\exp\!\left(-\frac{k}{2d4^d}\right).
\]
If $k\le(1+2d)2^d$, the asserted right side exceeds one, so the result is trivial.
The quotient argument in Corollary~\ref{cor:exponential-small-cuts} transfers this estimate to cyclic regions.
\end{proof}

\paragraph{Status and scope.}
The proof has passed independent mathematical audits, including an audit by Fable 5.2.
It is not formalized in Lean.
Its rate is $1/(2d4^d)$, instead of the earlier rate $2^{-O(d2^d)}$.
For an error margin $\varepsilon>0$, more than one half plus that margin forces
\[
 k<2d4^d\ln(2/\varepsilon).
\]
This improves the dependence on $\varepsilon$ from inverse-linear to logarithmic.

% New refinement. Root and path_stability independently audited all steps and constants.
\section{Conditioning on terminal data block by block}
\label{sec:second-order-block-conditioning}

A block is a maximal two-connected subgraph, with each bridge counted as a block.
An isolated vertex is a degenerate block when needed.
The cycle space of a graph is the direct sum of its block cycle spaces.
The block-cut graph of each connected component is a tree.
Its nodes are the blocks and the articulation vertices, with incidence as adjacency.

\begin{lemma}[Weighted conditional block estimate]\label{lem:weighted-block-conditioning}
Let $Q$ be a connected loopless multigraph, and let $S$ be a nonempty proper vertex set.
For each vertex $v$, let $L_v$ be its local edge-coordinate space in $\mathcal Z(Q)^*$.
Put
\[
 U=\sum_{s\in S}L_s,\qquad B=|\delta_Q(S)|.
\]
Let $\mathcal V\subseteq V(Q)\setminus S$ consist of vertices with $L_v\cap U\ne0$.
Write $I_v$ for the indicator that $v$ is closed, and put
$p_v=\mathbb P(I_v=1)=2^{-\dim L_v}$.
For arbitrary $0\le\eta_v\le1$, there is a set $E\subseteq\mathcal V$ with $|E|\le5B$ such that
\begin{equation}\label{eq:weighted-block-conditioning}
 \mathbb E\prod_{v\in\mathcal V}(1-\eta_v I_v)
 \le\frac12+\frac12\exp\!\left(-\sum_{v\in\mathcal V\setminus E}\eta_v p_v\right).
\end{equation}
\end{lemma}

\begin{proof}
Put $H=Q-S$.
In its block-cut forest, attach one distinct terminal leaf for each edge of $\delta_Q(S)$.
Attach that leaf to the articulation node of its endpoint in $H$, if that endpoint is an articulation.
Otherwise attach it to the unique block containing its endpoint.
Thus terminals at nonarticulation vertices are also retained.
Prune each tree to the union of paths between terminal leaves.
Components with fewer than two terminals contribute nothing.
Call the resulting forest $T$.
It has at most $B$ leaves.

Exclude the vertices of $\mathcal V$ adjacent to $S$, of which there are at most $B$.
The cut-space argument in Lemma~\ref{lem:local-terminal-intersection} applies to every remaining vertex.
It shows that such a vertex is an articulation separating at least two terminal-bearing components of $H-v$.
Its articulation node therefore has degree at least two in $T$.

Also exclude articulation nodes of degree at least three in $T$.
There are at most $B$ of these.
Finally exclude articulation nodes adjacent in $T$ to a block node of degree at least three.
The sum of the degrees of all branching nodes in a forest with at most $B$ leaves is at most $3B$.
This last exclusion therefore removes at most $3B$ further vertices.
Let $E$ be the union of the excluded vertices.

Each remaining vertex $v$ has exactly two terminal-bearing incident blocks in $H$.
All other incident blocks lie in branches containing no terminal.
The two terminal sides give the same nonzero functional $a_v\in U$.
The cut-space argument also shows
\begin{equation}\label{eq:block-intersection-line}
 L_v\cap U=\{0,a_v\}.
\end{equation}
Indeed, a local cut can only select a union of the components of $H-v$.
Terminal-free components contribute zero.
The two terminal-bearing components have complementary terminal sets, whose total functional is zero.

Join two remaining vertices when they belong to a common block of $H$.
Any shared block lies on the terminal forest between their articulation nodes.
It has degree two there, since neighbors of branching block nodes were excluded.
The resulting conflict graph is therefore a forest of paths.
Its two color classes contain no pair of vertices incident with the same block.
Choose a color class $J$ with
\[
 \sum_{v\in J}\eta_vp_v
 \ge\frac12\sum_{v\in\mathcal V\setminus E}\eta_vp_v.
\]

Condition on the evaluation of all functionals in $U$.
This determines every edge incident with $S$.
The remaining edge set in $H$ is a uniform member of an affine translate of $\mathcal Z(H)$.
Its random choices in different blocks are independent.

For clarity, zero-star feasibility can also be checked within each block.
At an articulation $v$, each incident block has a prescribed star parity, determined by the terminal data on its side.
The terminal-free sides have parity zero.
The other two sides have parity $a_v$.
Deleting $v$ from a nondegenerate block leaves it connected.
Hence, whenever the prescribed parity at $v$ is zero, a compatible edge set avoiding its star exists in that block.
This follows by joining its remaining odd vertices inside the connected graph obtained after deletion.
For a bridge block, parity zero simply forces its edge to be absent.
Thus the full star at $v$ can be zero exactly when $a_v=0$.

Equation~\eqref{eq:block-intersection-line} gives the exact conditional probability
\[
 \mathbb P(I_v=1\mid U)=2p_v\,\mathbf1_{\{a_v=0\}}.
\]
The indicators for $v\in J$ are conditionally independent, because their incident block sets are disjoint.
Ignoring the other factors and using $1-x\le e^{-x}$ gives
\[
 \mathbb E\!\left[\prod_{v\in\mathcal V}(1-\eta_vI_v)\,\middle|\,U\right]
 \le\exp\!\left(-2\sum_{v\in J}\eta_vp_v\mathbf1_{\{a_v=0\}}\right).
\]

Put $W=\sum_{v\in J}\eta_vp_v\mathbf1_{\{a_v=0\}}$ and $A=\sum_{v\in J}\eta_vp_v$.
Each $a_v$ is a nonzero functional, so $\mathbb EW=A/2$, while $0\le W\le A$.
If $A=0$, the claimed bound is trivial.
Otherwise the chord bound for the convex function $e^{-2x}$ gives
\[
 \mathbb E e^{-2W}\le\frac12+\frac12e^{-2A}.
\]
The choice of $J$ now proves~\eqref{eq:weighted-block-conditioning}.
No independence among the functionals $a_v$ is required.
\end{proof}

\begin{theorem}[A sharper exponential cut constant]\label{thm:sharper-exponential-cuts}
Let $Q$ be a connected loopless multigraph with $k\ge1$ marked vertices of degree at most $d\ge2$.
For a uniform even subgraph, let $Z$ count the closed marked vertices.
Then
\begin{equation}\label{eq:sharper-exponential-cuts}
 \mathbb E2^{-Z}\le\frac12+2\exp\!\left(-\frac{k}{8d2^d}\right).
\end{equation}
The same bound holds for simultaneous forest restrictions to $k$ disjoint connected cyclic regions with boundaries of size at most $d$.
\end{theorem}

\begin{proof}
A zero local space gives an expectation at most one half.
A direct sum of $2^d$ local spaces gives an expectation below one half, as before.
Otherwise choose a maximal direct-sum family with vertex set $S$ and size $t<2^d$.
Every other marked local space meets $U=\sum_{s\in S}L_s$ nontrivially.

Suppose first that $k\ge2(5d+1)2^d$.
Then $S$ is nonempty and proper.
Its boundary has size $B\le dt$.
Apply Lemma~\ref{lem:weighted-block-conditioning} with $\eta_v=1/2$ to the marked vertices outside $S$.
At least $k-t-5B$ marked vertices remain outside its exceptional set.
Each has $p_v\ge2^{1-d}$, so
\[
 \mathbb E2^{-Z}-\frac12
 \le\frac12\exp\!\left(-2^{-d}(k-t-5B)\right)
 \le\frac12\exp\!\left((5d+1)-\frac{k}{2^d}\right)
 \le\frac12\exp\!\left(-\frac{k}{2^{d+1}}\right).
\]
This is stronger than~\eqref{eq:sharper-exponential-cuts}.

If $k<2(5d+1)2^d$, then
\[
 \frac{k}{8d2^d}<\frac{5d+1}{4d}\le\frac{11}{8}<\ln4.
\]
The right side of~\eqref{eq:sharper-exponential-cuts} therefore exceeds one.
This proves the bound for all $k$.
The cyclic-region statement follows from the same quotient and internal-cycle-space argument as before.
\end{proof}

\paragraph{Weighted scope.}
Lemma~\ref{lem:weighted-block-conditioning} retains the exact closure probabilities $p_v$.
It loses only the weighted mass on its exceptional vertices and one half of the remaining mass in the path coloring.
Its bound depends on no maximum degree after the exceptional set has been chosen.
This identifies a possible route to a weighted estimate for regions with unequal boundaries.

% Weighted refinement; path_stability independently audited all steps and constants.
\section{A closure-mass version of the cut estimate}
\label{sec:closure-mass-cuts}

The degree bound and the local coordinate rank play different roles.
Keeping them separate gives a weighted bound with polynomial dependence on the maximum boundary size.

\begin{lemma}[Degree and local rank separately]\label{lem:degree-rank-cut-estimate}
Let $Q$ be connected and loopless, with $k$ distinct marked vertices.
Suppose their degrees are at most $d\ge2$ and their local coordinate spaces have dimension at most $s\ge1$.
If $Z$ counts the closed marked vertices in a uniform even subgraph, then
\[
 \mathbb E2^{-Z}\le\frac12+2\exp\!\left(-\frac{k}{8d2^{s+1}}\right).
\]
\end{lemma}

\begin{proof}
A zero local space gives an expectation at most one half.
Otherwise put $T=2^{s+1}$.
If $T$ marked local spaces form a direct sum, their closure indicators are independent.
Each has closure probability at least $2^{-s}$, so the expectation is at most
$(1-2^{-s-1})^T\le e^{-1}<1/2$.

Otherwise choose a maximal direct-sum family $S$ with $t<T$ vertices.
Every other marked local space meets $U=\sum_{v\in S}L_v$ nontrivially.
For $k\ge2(5d+1)T$, the set $S$ is nonempty and proper, and its boundary has size $B\le dt$.
Lemma~\ref{lem:weighted-block-conditioning}, with $\eta_v=1/2$, gives
\[
 \mathbb E2^{-Z}-\frac12
 \le\frac12\exp\!\left(-2^{-s-1}(k-t-5B)\right)
 \le\frac12\exp\!\left((5d+1)-\frac{k}{T}\right)
 \le\frac12 e^{-k/(2T)}.
\]
This is stronger than the stated bound.
For $k<2(5d+1)T$, one has $k/(8dT)<11/8<\ln4$.
The stated right side then exceeds one.
\end{proof}

\begin{corollary}[Total closure mass]\label{cor:closure-mass-cut-estimate}
Let $Q$ be connected and loopless, with distinct marked vertices of degree at most $d\ge2$.
For a uniform even subgraph, put
\[
 p_v=\Pr(v\text{ is closed}),\qquad \mu=\sum_{v\text{ marked}}p_v.
\]
Then
\begin{equation}\label{eq:closure-mass-cut-estimate}
 \mathbb E2^{-Z}
 \le\frac12+2\exp\!\left(-\frac{\mu}{16d(d-1)}\right).
\end{equation}
The same estimate bounds simultaneous forest restrictions in disjoint connected cyclic regions,
where $p_v$ is the probability that the corresponding region boundary is closed.
\end{corollary}

\begin{proof}
If some local space has dimension zero, the expectation is already at most one half.
Otherwise each dimension belongs to $\{1,\ldots,d-1\}$.
For each $s$, let $k_s$ count marked vertices of local dimension exactly $s$.
Their closure mass is $\mu_s=k_s2^{-s}$.
Some $s$ satisfies $\mu_s\ge\mu/(d-1)$.
Ignore the other marked factors and apply Lemma~\ref{lem:degree-rank-cut-estimate} to this class.
Its exponent is
\[
 \frac{k_s}{8d2^{s+1}}=\frac{\mu_s}{16d}
 \ge\frac{\mu}{16d(d-1)}.
\]
This proves~\eqref{eq:closure-mass-cut-estimate}.
The cyclic-region statement follows from the quotient and internal-cycle-space reduction.
\end{proof}

This bound retains the actual closure mass, which can be much larger than the estimate obtained from the maximum degree alone.
The regions must still be vertex-disjoint.
The estimate does not yet control an overlapping family of moderate-length cycles.

\section{The degree dependence in exponential small-cut decay}
\label{sec:cut-rate-obstruction}

The factor $2^{-d}$ in the improved exponential rate is necessary up to a polynomial factor in $d$.
Parallel bundles on a cycle give an explicit example.

Let $Q_{k,m}$ replace every edge of a $k$-cycle by $m$ parallel edges, where $k\ge3$ and $m\ge1$.
Mark all $k$ vertices, each of degree $d=2m$.
For a uniform even subgraph, let $Z$ count the closed marked vertices.
Put $s=2^{1-m}$ and define
\[
 T_s=\begin{pmatrix}1-s&s\\1-s&s/2\end{pmatrix},
 \qquad
 \lambda_\pm(s)=\frac{1-s/2\pm\sqrt{1+s-7s^2/4}}2.
\]

\begin{proposition}\label{prop:cut-rate-obstruction}
For these graphs,
\begin{equation}\label{eq:bundle-cycle-moment}
 \mathbb E2^{-Z}
 =\frac12+\frac12\operatorname{tr}(T_s^k)
 =\frac12+\frac12\bigl(\lambda_+(s)^k+\lambda_-(s)^k\bigr).
\end{equation}
For fixed $m$, the exponential decay rate of the excess above one half is $-\ln\lambda_+(s)$.
As $m\to\infty$,
\[
 -\ln\lambda_+(s)=\frac{s^2}{2}+O(s^3)
 =(1+o(1))2^{1-d}.
\]
Consequently, any bound of the form
\[
 \mathbb E2^{-Z}\le\frac12+C_d e^{-c_d k}
\]
with finite $C_d$ independent of $k$ must satisfy $c_d\le-\ln\lambda_+(2^{1-d/2})$ for even $d$.
\end{proposition}

\begin{proof}
Parity at each vertex forces all bundle parities to agree.
Their common parity is a uniform bit.
Conditional on either value, the edge choices within different bundles are independent and uniform among the $2^{m-1}$ choices of that parity.

In the odd sector every bundle is nonempty, so every marked vertex is open.
In the even sector, a bundle is empty with probability $s$.
Let $X_1,\ldots,X_k$ be the independent indicators of empty bundles in this sector.
Then
\[
 2^{-Z}=\prod_{i=1}^k\left(1-\frac12X_iX_{i+1}\right),
 \qquad X_{k+1}=X_1.
\]
Summing over the Bernoulli variables gives the transfer-matrix trace $\operatorname{tr}(T_s^k)$.
The entries are $T_s(a,b)=\mathbb P(X=b)(1-ab/2)$ for $a,b\in\{0,1\}$.
This proves~\eqref{eq:bundle-cycle-moment}.

The trace and determinant of $T_s$ are $1-s/2$ and $-s(1-s)/2$.
They give the displayed eigenvalues.
For $0<s<1$, one has $\lambda_+>0$ and $|\lambda_-|<\lambda_+$.
For $s=1$, the eigenvalues are $1/2$ and zero.
Thus the dominant eigenvalue gives the stated fixed-$m$ decay rate.
Expanding its explicit formula at $s=0$ gives $\lambda_+(s)=1-s^2/2+O(s^3)$.
The last assertion follows by letting $k\to\infty$ with $d$ fixed.
\end{proof}

The exact closure probability at each vertex is $p=2^{1-d}=s^2/2$.
Thus this family has total closure mass $\mu=kp$ and excess decay rate $(1+o(1))\mu$ as $d$ grows.
It supports a possible weighted formulation in terms of closure mass.
It rules out replacing the degree dependence by a substantially larger exponential rate uniformly over these multigraphs.


\section{Verification and scope}
The mathematical argument applies to every finite connected graph satisfying its hypotheses.
The finite checks test the structural identities and conditional laws used in the proof.
They are not substitutes for those proofs.

The script \nolinkurl{check_exponential_small_cuts.py} uses exact binary linear algebra and rational probabilities.
Its initial catalogue contains 108 connected simple graphs of maximum degree four through order six and 16 additional multigraphs.
It checks terminal partitions, common-functional splits, and full conditional laws.
The saved record includes 5,682 terminal subsets, 878 common-functional splits, and 1,028 nontrivial block-conditioning checks.

The script \nolinkurl{check_bundle_cycle_decay.py} compares three independent exact counts.
These are cycle-space enumeration, sums over empty-bundle indicators, and transfer-matrix traces.
Its 12 cases contain 77,558 cycle-space states.
All comparisons agree.

The bounded regions must be vertex-disjoint.
The present result does not directly control an overlapping family of moderate-length cycles.
That distinction remains important for the unrestricted second-order cycle bound.
The existing main manuscripts, formal sources, and earlier reviewed reports remain unchanged.
\end{document}
